% !TEX root = ../main.tex
% ================== MOTIVACAO + INTRODUCAO ==================
\secao{Motivação}

\begin{itemize}
  \item Sistemas de controle de temperatura em malha fechada são a base de fornos,
        incubadoras, estufas e diversos processos industriais, mas plataformas didáticas
        para estudá-los costumam ser caras, volumosas ou de difícil instrumentação.
  \item Propor uma bancada compacta, de baixo custo e totalmente instrumentada para observar
        e validar, de forma visual e mensurável, um controlador PI atuando sobre um sistema
        térmico real.
\end{itemize}

\secao{Introdução}

O controle de temperatura em malha fechada combina instrumentação, eletrônica de potência e
teoria de controle. Este trabalho apresenta uma câmara térmica compacta
(\(\approx\)100\,×\,100\,×\,80\,mm) aquecida por três lâmpadas automotivas W5W, com sensor
LM35 e um controlador PI embarcado em um ESP32, que ajusta o PWM aplicado às lâmpadas por meio
de um MOSFET IRLZ44N para manter a temperatura interna em um setpoint definido [1, 5].

\titulofig{Objetivos específicos}
\begin{itemize}
  \item Estudar controle PI, PWM e instrumentação com sensores analógicos;
  \item Projetar e montar a câmara (acrílico, lâmpadas W5W, chapa de alumínio, LM35 e cadeia
        de potência protegida);
  \item Implementar o firmware do ESP32 (ADC, PWM via LEDC, PI com anti-windup e compensação
        de bateria);
  \item Realizar o ensaio de degrau, sintonizar o PI e validar a malha fechada.
\end{itemize}
